\section{Experimental Settings}

\subsection{Research Questions}
The main research question that this project aims to address is:
\begin{itemize}
    \item \textbf{RQ1}: Can the integration of aligned multimodal embeddings (text, image, audio) enhance the performance of recommender systems compared to unimodal or non-aligned multimodal approaches?
\end{itemize}

% =============================================================================
\subsection{Datasets}

This project employs two benchmark datasets for evaluating multimodal recommendation:

\subsubsection{MovieLens 1M}

MovieLens 1M is a widely-used collection of explicit movie ratings containing \textbf{User-movie interactions}: implicit feedback from user ratings.
For each movie, the dataset has been enriched with multimodal data:
\begin{itemize}
    \item \textbf{Images}: movie posters (.jpg format)
    \item \textbf{Audio}: movie trailer (.wav format)
    \item \textbf{Text}: textual descriptions, synopses, or reviews (.txt format)
\end{itemize}

\subsubsection{LastFM}

LastFM is a music listening dataset containing:
\begin{itemize}
    \item \textbf{User-artist interactions}: implicit feedback from listening history
    \item \textbf{Multiple tracks per artist}: track-level multimodal features
\end{itemize}

For each track, the dataset includes:
\begin{itemize}
    \item \textbf{Images}: album artwork (.jpg format)
    \item \textbf{Audio}: songs (.wav format)
    \item \textbf{Text}: track metadata or descriptions (.txt format)
\end{itemize}

\textbf{Note}: LastFM presents a unique challenge where interactions are recorded at the artist level, while multimodal features are extracted at the track level. This mismatch is addressed in the mapping pipeline (Section~\ref{sec:mapping}).

% =============================================================================
\subsection{Multimodal Embedding Extraction Pipeline}
\label{sec:extraction}

The embedding extraction pipeline is \textbf{dataset-agnostic} and applies the same process to both MovieLens and LastFM. This section describes the unified extraction methodology.

\subsubsection{Overview}

The extraction pipeline performs the following steps:
\begin{enumerate}
    \item File discovery and alignment across modalities
    \item Validation and corruption detection
    \item Per-modality embedding extraction using multiple encoder architectures
    \item Strict consistency enforcement across modalities
    \item Output generation in aligned NumPy arrays
\end{enumerate}

\subsubsection{File Organization and Naming Convention}

All multimodal files are organized with a consistent basename structure. For both datasets:
\begin{itemize}
    \item Each item (movie/track) is identified by a unique ID
    \item Files across modalities share the same basename
\end{itemize}

\textbf{Example (MovieLens)}: For movie ID \texttt{1}:
\begin{itemize}
    \item Image: \texttt{ml1m/\_images/1.jpg}
    \item Audio: \texttt{ml1m/\_audios/1.wav}
    \item Text: \texttt{ml1m/\_texts/1.txt}
\end{itemize}

\textbf{Example (LastFM)}: For artist \texttt{1000} with track variant \texttt{1}:
\begin{itemize}
    \item Image: \texttt{lastfm/\_images/1000\_1.jpg}
    \item Audio: \texttt{lastfm/\_audios/1000\_1.wav}
    \item Text: \texttt{lastfm/\_texts/1000\_1.txt}
\end{itemize}
(there are up to 5 songs per artist)

\subsubsection{Embedding Models}

Multiple encoder architectures are employed to extract embeddings, providing both aligned and specialized representations:

\paragraph{Aligned Multimodal Models}
\begin{itemize}
    \item \textbf{AudioCLIP}: Extracts aligned embeddings across all three modalities (image, audio, text) in a shared feature space
    \item \textbf{CLIP}: Extracts aligned image-text embeddings
\end{itemize}

\paragraph{Modality-Specific Models}
\begin{itemize}
    \item \textbf{ViT (Vision Transformer)}: Image embeddings
    \item \textbf{VGGish}: Audio embeddings
    \item \textbf{MiniLM}: Text embeddings
\end{itemize}



\subsubsection{Extraction Process}

\paragraph{Step 1: File Discovery and Alignment}

Directories containing raw multimodal data are scanned:
\begin{itemize}
    \item Identifies files by basename (e.g., movie ID or artist\_track combination)
    \item Computes the \textbf{intersection} of basenames across all three modalities
\end{itemize}

Only items present in all three modalities are considered candidates for extraction. This ensures row-wise correspondence in the final embedding arrays.


\paragraph{Step 2: Extraction}
Files that pass validation are processed in order to extract embeddings. If the extraction for a specific modality fails (e.g., due to file corruption), the entire item is excluded to ensure strict alignment.



% =============================================================================
\subsection{Interaction Filtering and Remapping Pipeline}
\label{sec:mapping}

After embedding extraction, the interaction datasets must be filtered and remapped to ensure consistency with the available multimodal features. This section describes the \textbf{dataset-specific} mapping procedures.

\subsubsection{Common Pipeline Steps}

Both datasets follow the same core logic:
\begin{enumerate}
    \item Load original interaction file and a .csv mapping of item IDs to embedding row indices
    \item Apply 5-core filtering to ensure interaction density
    \item Remap user and item IDs to contiguous ranges [0, ..., n-1]
    \item Reconstruct embedding arrays in the new item order
    \item Generate output files with remapped IDs and embeddings
\end{enumerate}

The key difference lies in \textbf{Step 2}, where LastFM requires special handling for artist-to-track mapping.

\subsubsection{MovieLens Mapping}

\paragraph{Step 1: Loading Data}
The script loads:
\begin{itemize}
    \item \texttt{movielens\_1m.inter}: Original interaction file (user\_id, item\_id, rating, timestamp)
    \item A .csv mapping from movie ID to embedding row index
    \item All \texttt{.npy} embedding files
\end{itemize}

\paragraph{Step 2: Feature-Based Filtering}
Check between \texttt{movielens\_1m.inter} and the .csv file to retain only interactions referencing movies contained in both files.

\paragraph{Step 3: 5-Core Filtering}
To address data sparsity, the script applies iterative 5-core filtering:

\textbf{Algorithm}:
\begin{enumerate}
    \item Count interactions per user and per item
    \item Remove users with $< 5$ interactions
    \item Remove items with $< 5$ interactions
    \item Repeat until convergence (no further removals)
\end{enumerate}

This ensures that every user has rated at least 5 items and every item has been rated by at least 5 users.

\paragraph{Step 4: ID Remapping}
After filtering, IDs may be non-contiguous. The script creates contiguous mappings:
\begin{itemize}
    \item \textbf{Users}: $\{u_1, u_2, ..., u_m\} \rightarrow \{0, 1, ..., m-1\}$
    \item \textbf{Items}: $\{i_1, i_2, ..., i_n\} \rightarrow \{0, 1, ..., n-1\}$
\end{itemize}

\paragraph{Step 5: Embedding Reindexing}
After remapping, the embedding arrays are reconstructed to match the new item order and a new .csv mapping file is generated.


\subsubsection{LastFM Mapping}

LastFM requires special handling due to a structural mismatch:
\begin{itemize}
    \item \textbf{Interactions} are recorded at the \textbf{artist level} (user listens to artist)
    \item \textbf{Embeddings} are extracted at the \textbf{track level} (one feature set per track)
\end{itemize}

\paragraph{Resolution Strategy: Artist-to-Track Expansion}

The mapper resolves this mismatch by:
\begin{enumerate}
    \item Mapping each artist to all available track variants in the .csv file
    \item Expanding each user-artist interaction into multiple user-track interactions
    \item Treating each variant as a separate item in the recommendation task
\end{enumerate}



\paragraph{Step 1: Loading Data}
The script loads:
\begin{itemize}
    \item \texttt{lastfm.inter}: Original interaction file (user\_id, artist\_id, listen\_count, timestamp)
    \item A .csv mapping from \texttt{artist\_track} identifier to embedding row index
    \item All \texttt{.npy} embedding files
\end{itemize}

\paragraph{Step 2: Artist-to-Track Mapping and Expansion}
The script performs two sub-steps:

\textbf{Sub-step 2a: Build Artist-to-Tracks Dictionary}
\begin{itemize}
    \item Parse the .csv file to extract artist and track identifiers
    \item Group tracks by artist: $\text{artist} \rightarrow [\text{track}_1, \text{track}_2, ...]$
\end{itemize}

\textbf{Sub-step 2b: Expand Interactions}
For each user-artist interaction :
\begin{itemize}
    \item Lookup available tracks for artist $a$: $T_a = [\text{track}_1, ..., \text{track}_k]$
    \item Create $k$ new interactions: $(u, \text{track}_1), ..., (u, \text{track}_k)$
\end{itemize}

This expansion significantly increases the number of interactions.

\paragraph{Step 3--5: Standard Pipeline}

After expansion, the pipeline follows the same steps as MovieLens

% =============================================================================
\subsection{Data Distribution Extraction}

After extracting and remapping embeddings, it is essential to analyze their geometric and statistical properties to understand how different encoders structure the feature space.

\subsubsection{Overview}

The distribution extraction pipeline performs the following steps:
\begin{enumerate}
    \item Load pre-extracted embedding arrays (\texttt{.npy} files)
    \item Apply UMAP for dimensionality reduction to 2D
    \item Normalize reduced embeddings onto the unit circle (S1 manifold)
    \item Compute 2D Gaussian KDE for spatial density visualization
    \item Compute angular KDE for directional distribution analysis
\end{enumerate}

\subsubsection{UMAP Dimensionality Reduction}

High-dimensional embeddings are reduced to 2D using UMAP (Uniform Manifold Approximation and Projection).

\paragraph{Algorithm Overview}

UMAP is technique used to projects an high-dimensional space into a lower-dimensional soace while preserving both local and global structure. The algorithm operates in three main phases:

\begin{enumerate}
    \item \textbf{Graph Construction}: Build a weighted k-nearest neighbor graph in high-dimensional space, where edge weights represent fuzzy membership strengths based on local distance distributions
    \item \textbf{Fuzzy Simplicial Set}: Convert the neighbor graph into a fuzzy topological structure that captures the geometry
    \item \textbf{Optimization}: Minimize the cross-entropy between the high-dimensional fuzzy set and its low-dimensional representation using stochastic gradient descent
\end{enumerate}


\paragraph{S1 Normalization}

After UMAP reduction, all points are projected onto the unit circle:
\begin{enumerate}
    \item Center the data: $\mathbf{x}_i' = \mathbf{x}_i - \bar{\mathbf{x}}$
    \item Normalize by Euclidean norm: $\mathbf{x}_i'' = \frac{\mathbf{x}_i'}{||\mathbf{x}_i'||}$
\end{enumerate}


\subsubsection{Density Estimation}

Two complementary density estimation methods are applied to quantify the spatial and directional distribution of embeddings.

\paragraph{ 2D Gaussia Kernel Density Estimation (KDE) Overview}

KDE is a non-parametric method for estimating the probability density function of a random variable. Given a set of data points $\{\mathbf{x}_1, \ldots, \mathbf{x}_n\}$, the kernel density estimate at position $\mathbf{x}$ is:

\[
\hat{f}(\mathbf{x}) = \frac{1}{n} \sum_{i=1}^{n} K_h(\mathbf{x} - \mathbf{x}_i)
\]

where $K_h$ is a kernel function (in this case the Gaussian kernel) 


This reveals regions of high and low embedding density, indicating clustering patterns and feature space coverage.

\paragraph{Angular KDE}

The distribution of angular positions is analyzed separately:
\begin{enumerate}
    \item Convert Cartesian coordinates to polar angles: $\theta_i = \arctan2(y_i, x_i)$
    \item Fit Gaussian KDE on angular values: $\theta \in [-\pi, \pi]$
    \item Plot density function as a 1D curve
\end{enumerate}

This analysis reveals if embeddings are uniformly distributed or not




